% from ega0-23-fr.tex
\oldpage[0\textsubscript{IV-1}]{214}

un $A$-module de type fini (autrement dit, une $A$-algèbre finie). On dit
qu'un anneau $A$ est universellement japonais si toute $A$-algèbre intègre de
type fini est un anneau japonais.
\end{definition}

Il est clair que si $A$ est un anneau japonais, tout anneau de fractions
$S^{-1}A$ est alors un anneau japonais, car (avec les notations précédentes)
$S^{-1}A'$ est la fermeture intégrale de $S^{-1}A$ dans $K'$, et est
évidemment un $S^{-1}A$-module de type fini.

Un anneau intègre noethérien, même un \emph{anneau de valuation discrète},
n'est pas nécessairement un anneau japonais [30].

\begin{proposition}[23.1.2]
\label{0.23.1.2-fr}
Soient $A$ un anneau noethérien intègre, $K$ son corps des fractions. Si, pour
toute extension radicielle finie $K'$ de $K$, la fermeture intégrale de $A$
dans $K'$ est un $A$-module de type fini, alors $A$ est un anneau japonais.
\end{proposition}

\begin{proof}
Comme $A$ est noethérien, il suffit, pour vérifier que $A$ est un anneau
japonais, de voir que pour toute extension quasi-galoisienne finie $L$ de
$K$, la fermeture intégrale $B$ de $A$ dans $L$ est un $A$-module de type
fini. Or $L$ est une extension galoisienne de la plus grande extension
radicielle $K'$ de $K$ contenue dans $L$~; et si $A'$ est la fermeture
intégrale de $A$ dans $K'$, $B$ est la fermeture intégrale de $A'$ dans $L$.
Mais on sait que dans une extension séparable de $K'$, la fermeture intégrale
de $A'$ est un $A'$-module de type fini (Bourbaki, \emph{Alg. comm.},
chap.~V, §~1, n\textsuperscript{o}~6, cor.~1 de la prop.~20), d'où la
proposition.
\end{proof}

Il résulte de (23.1.2) que lorsque $K$ est de caractéristique $0$, dire que
$A$ est un anneau japonais signifie que sa \emph{clôture intégrale} $A'$ est
un $A$-module de type fini.

\begin{theorem}[23.1.3]
\label{0.23.1.3-fr}
\textnormal{(Tate).} Soient $A$ un anneau intègre noethérien, $x\neq0$ un
élément de $A$. Supposons vérifiées les conditions suivantes~:
\begin{enumerate}
  \item[(i)] $A$ est intégralement clos.
  \item[(ii)] $\mathfrak p=xA$ est premier et $A$ est séparé et complet pour
  la topologie $\mathfrak p$-préadique.
  \item[(iii)] $A/xA$ est un anneau japonais.
\end{enumerate}
Alors $A$ est un anneau japonais.
\end{theorem}

Nous utiliserons le lemme suivant~:

\begin{lemma}[23.1.3.1]
\label{0.23.1.3.1-fr}
Soient $A$ un anneau, $x$ un élément de $A$ non diviseur de $0$ dans $A$ et
tel que $xA=\mathfrak p$ soit premier~; alors, pour tout entier $n>0$,
l'image réciproque de $x^nA_{\mathfrak p}$ dans $A$ est $x^nA$.
\end{lemma}

\begin{proof}
En effet, supposons que $b\in A$ soit un élément tel que
$b/1=x^na/s$ dans $A_{\mathfrak p}$, où $a\in A$ et
$s\notin\mathfrak p$~; il existe donc $s'\notin\mathfrak p$ tel que
$s'sb=x^nas'$, d'où $s'sb\in\mathfrak p$, et comme
$s's\notin\mathfrak p$, cela entraîne $b\in\mathfrak p$, autrement dit
$b=xb'$ avec $b'\in A$~; puisque $x$ est non diviseur de $0$, on en conclut
$s'sb'=x^{n-1}as$ et il suffit de raisonner par récurrence sur $n$.
\end{proof}

\begin{proof}
Pour démontrer le théorème, on peut se borner au cas où le corps des fractions
$K$ de $A$ est de caractéristique $p>0$, puisque $A$ est intégralement clos.
Soit $K'$ une extension radicielle finie de $K$, de sorte qu'il existe une
puissance $q=p^e$ telle que $K'^q\subset K$~; en remplaçant $K'$ par une
extension radicielle plus grande, on peut même supposer qu'il existe $y\in K'$
tel que $y^q=x$. Soit $A'$ la fermeture intégrale de $A$ dans $K'$~; comme
$A$ est intégralement clos, $A'$ est l'ensemble des $x'\in K'$ tels que
$x'^q\in A'\cap K=A$. Posons $V=A_{\mathfrak p}$~; l'idéal maximal
$\mathfrak m=xV$ de $V$ étant principal, on sait que l'anneau local
noethérien $V$ est un \emph{anneau de valuation discrète} (17.1.4)~; comme
$V$ est intégralement clos, le même

